\section{Discussion, Scope, and the Handoff to Reproductive Accretion}
\label{sec:discussion}

Several limitations are deliberate.

First, the capability stock $q$ is reduced form. It bundles technical, organizational, labor, supplier, managerial, and market-facing capabilities that would need separate measurement in empirical work. The scalar state is retained because it makes the graduation mechanism transparent. The split into $(q^S,q^R)$ occurs only when the paper reaches exit, where transferability becomes unavoidable.

Second, the distinction between durable capability and contemporaneous dependence is conceptually clean but empirically difficult. Researchers must decide whether an observed complement created a transferable stock or merely provided a service that disappears with the relationship. The provenance object $\Delta q_e^u(T)$ is even more demanding because it requires a defensible no-support counterfactual. The present paper identifies the questions but does not solve their measurement.

Third, the activation model assumes nonnegative complementarity. That is useful for the entrance problem but cannot describe the whole economy. Competition for labor and capital, import competition, exchange-rate effects, congestion, and demand displacement are signed interactions. A fuller structural-change model would need something like $J=J^+-J^-$ or a general-equilibrium production-network environment. The paper therefore does not invoke Tarski-style monotonicity after signed dynamics enter.

Fourth, Proposition~\ref{prop:spectral-stability} is deliberately local. It holds an active configuration and external environment fixed and uses a linear capability-conversion law. The spectral boundary identifies whether that local linear system settles to a finite state; it does not imply that real capabilities are globally linear or that a supercritical linear specification should literally diverge forever. Nonlinear saturation is an open extension.

Fifth, the graduation frontier in equation~\eqref{eq:frontier} is a local possibility set, not a measured national development index. A country may have many relations with $\gamma>1$ that are never coordinated into existence. Conversely, a relation with $\gamma<1$ today may become graduable after infrastructure, external access, technology transfer, or the graduation of complements changes its environment.

Sixth, the model does not establish a welfare criterion for preserving strategic capability. Whether a country should retain a high-cost domestic industry because of reconstitution time, geopolitical risk, or supply-chain leverage is a separate problem. That question requires comparing the cost of maintaining capability with the expected value of resilience, substitution, and strategic options. It is therefore left to a subsequent paper rather than folded into the present one.

Seventh, the model does not make industrial persistence synonymous with human welfare. A regime could raise capability conversion by imposing extreme pressure on workers or students. The welfare layer must distinguish discipline applied to firms or supported positions from the viability and agency of people. A useful policy principle is ``hard discipline for firms; real transition possibilities for people,'' but that principle is not part of the positive theorem developed here.

The paper's natural endpoint is the production of an evolving set of graduated, generative, displaced, and recombined productive relations. A separate analysis of reproductive accretion asks a different question: when do such relations form a higher-order network capable of reproducing its productive state, and how robust is that closure to edge or node loss?

The bridge is no longer wholly formal terra incognita. Proposition~\ref{prop:spectral-stability} shows that this paper's endogenous capability dynamics already form a positive coupled system with a spectral boundary. The reproductive-accretion analysis studies a different substantive object---productive reinforcement and higher-order closure---but it belongs to the same mathematical family of nonnegative coupled systems. Conceptually, the handoff is
\begin{equation}
\begin{aligned}
\underbrace{
\text{activation}\rightarrow\text{capability}\rightarrow\text{graduation}
}_{\text{formation and graduation}}
&\rightarrow
\underbrace{
\text{capability amplification}\rightarrow\text{generativity}
}_{\text{partial bridge}}\\
&\rightarrow
\underbrace{
\text{accretion}\rightarrow\text{reproductive closure}\rightarrow\text{robustness}
}_{\text{separate analysis}}.
\end{aligned}
\label{eq:handoff}
\end{equation}
The substantive map remains open: one still needs to specify how the capability states and active relations in this paper generate the productive-reinforcement weights of the accretion model. A plausible bridge is a time-varying weight
\begin{equation}
w_{ef,t}=\phi_{ef}(q_t,x_t,\Omega_t),
\end{equation}
but specifying and validating $\phi$ is beyond the scope of the present draft. The claim is therefore modest: the two analyses already share a positive-systems architecture, while the economic mapping between their state variables remains to be identified.
